\chapter{Envelopes before the general theorem}
\index{envelope property!history}\index{unibranch variety}
\label{chap:bj-envelope-history}

An upper envelope is the nonlinear closure operation that makes
pluripotential theory work.  In the trivially valued setting its regularity
was the principal missing foundational statement in the theory developed by
Boucksom--Jonsson.  This chapter freezes the historical picture immediately
before the theorem proved in
\Crefrange{chap:bj-envelope-statement}{chap:bj-envelope-consequences}.  We
establish the formal equivalences that do not depend on
that theorem, state the cases known from Boucksom--Jonsson and their
addendum, and isolate exactly what remained open.

Throughout, $X$ is an integral projective variety over an algebraically
closed trivially valued field and $\theta\in N^1(X)$.  For a function
$f:X^{\An}\to\mathbf R\cup\{\pm\infty\}$, write
\[
 f^*(v):=\limsup_{w\to v}f(w)
\tag{36.1}\label{eq:bj-usc-regularization}
\]
for its upper-semicontinuous regularization.  The limit superior is
topological and is defined using neighbourhoods, not a chosen sequence.


\section{Orientation}
\label{sec:ch36-orientation}

\subsection{What is at stake}

Every construction in \Cref{chap:bj-fs-pl-psh} produced $\theta$-psh
functions by \emph{decreasing} approximation, and
\Cref{thm:bj-basic-closure} closed the class under generically finite
decreasing limits, finite maxima, and convex combinations.  It said nothing
about suprema.
That omission is not an oversight of the exposition: it is the precise shape
of the conjecture left by the published theory and its addendum, and proved in
\Crefrange{chap:bj-envelope-statement}{chap:bj-envelope-consequences}.

The reason suprema matter is that essentially every existence argument in
pluripotential theory is an envelope argument.  One writes down the family of
all psh competitors below an obstacle, takes the supremum, and needs to know
that the result is again psh.  Without that, the variational method has no
output, the projection $P_\theta$ is not an operator on the psh class, and
the metric geometry of finite energy has no maximal elements to work with.

\subsection{The three distinctions that organize this chapter}

A single imprecision has caused more confusion in this subject than any
other: treating the following three operations as interchangeable.

\begin{enumerate}
\item A \emph{decreasing net} $\phi_i\searrow\phi$.  The limit is
      automatically usc, and \Cref{thm:bj-basic-closure}\,(4) keeps it in the
      class provided it is generically finite; otherwise the closure theorem
      allows the identically $-\infty$ alternative.  This is how
      $\PSH_X(\theta)$ was defined.
\item An \emph{increasing net} $\phi_i\nearrow\phi$, or equivalently the
      supremum of a directed family.  The pointwise limit need not be usc,
      and the regularization $(\sup_i\phi_i)^*$ may differ from
      $\sup_i\phi_i$ on a nonempty set.  Whether the regularization is psh
      is the envelope property.
\item An arbitrary \emph{bounded-above family}, not directed.  Formally more
      general than~(2), but in fact equivalent to it, because the finite
      maxima of a family form a directed set with the same supremum and
      $\PSH_X(\theta)$ is closed under finite maxima.  This reduction is
      the content of \Cref{thm:bj-envelope-equivalences} and is the one
      simplification in this circle of ideas that is free.
\end{enumerate}

A fourth temptation is to replace every net by a sequence.  The formal
reduction of an arbitrary family to its finite maxima does not permit this:
its directed set need not have a countable cofinal subset.  This precise
obstruction is recorded in \Cref{sec:ch36-nets}; it does not assert that a
particular psh function can never admit a sequential approximation.

\subsection{What this chapter proves, and what it does not}

Proved here, independently of the envelope theorem: the formal calculus of
the obstacle envelope $P_\theta$; the equivalence of the envelope
property with normalized compactness and with an obstacle criterion; the
implication that the envelope property in an ample class forces
unibranchness, so that the hypothesis of the main theorem is not an artefact
of its proof; and the two published positive results for smooth varieties,
first for nef classes and then, in the addendum, for arbitrary classes under
the same characteristic--dimension hypotheses.  The
addendum also proves a separate energy identity on possibly singular
varieties under the weak envelope property.

Not proved here, and not used: the envelope theorem itself.  The firewall is
stated precisely at the end of the chapter, because a circular dependency in
exactly this place would be easy to introduce and hard to detect.

\section{Two envelope operations}

\begin{definition}[Envelope property]
\label{def:bj-envelope-property}
The class $\theta$ has the envelope property, denoted
$\operatorname{EP}_X(\theta)$, if every nonempty bounded-above family
$(\phi_i)_{i\in I}\subset\PSH_X(\theta)$ satisfies
\[
                        \left(\sup_{i\in I}\phi_i\right)^*
                              \in\PSH_X(\theta).
\tag{36.2}\label{eq:bj-envelope-property}
\]
\end{definition}

Bounded above means bounded by one real constant on all of $X^{\An}$.
It does not mean that each $\phi_i$ is bounded below.  In particular, the
members may equal $-\infty$ on proper-support points.

\begin{definition}[Psh obstacle envelope]
\label{def:bj-obstacle-envelope}
For $h:X^{\An}\to\mathbf R\cup\{\pm\infty\}$, define
\[
 P_\theta(h):=
 \sup\{\psi\in\PSH_X(\theta)\mid\psi\leq h\}.
\tag{36.3}\label{eq:bj-obstacle-envelope}
\]
If there is no competitor, $P_\theta(h)\equiv-\infty$.
\end{definition}

Despite its name, the raw supremum in \eqref{eq:bj-obstacle-envelope} need
not be upper semicontinuous or psh without an envelope theorem.  It may also
be identically $+\infty$ when the competitors are not uniformly bounded
above.  These are three distinct alternatives.

\begin{lemma}\label{lem:obstacle-envelope-formal}
For arbitrary obstacles $h,g$ and $c\in\mathbf R$:
\begin{enumerate}
\item $h\leq g$ implies $P_\theta(h)\leq P_\theta(g)$;
\item $P_\theta(h+c)=P_\theta(h)+c$;
\item $P_\theta(P_\theta(h))=P_\theta(h)$ whenever the inner expression
      is used as an obstacle;
\item if $h_j\uparrow h$ is an increasing net of bounded-below lower
      semicontinuous functions, then
      $P_\theta(h_j)\uparrow P_\theta(h)$.
\end{enumerate}
\par
\end{lemma}

\begin{proof}
The first three claims follow by comparing the defining competitor sets.
For (4), monotonicity gives
$\sup_jP_\theta(h_j)\leq P_\theta(h)$.  Let
$\psi\in\PSH_X(\theta)$ satisfy $\psi\leq h$, and fix
$\varepsilon>0$.  Since $\psi$ is usc, $h_j$ is lsc, and
$X^{\An}$ is compact, the increasing-net version of Dini's covering
argument gives $\psi<h_j+\varepsilon$ for all sufficiently large $j$.
Thus $\psi-\varepsilon\leq P_\theta(h_j)$.  Taking the supremum over
$\psi$, then letting $\varepsilon\downarrow0$, proves the reverse
inequality.
\end{proof}

\section{Increasing nets and normalized compactness}

The finite maximum of two $\theta$-psh functions is $\theta$-psh.
Hence any family can be replaced by the increasing net of its finite maxima,
indexed by the directed set of finite subsets.

\begin{theorem}\label{thm:bj-envelope-equivalences}
For any $\theta\in N^1(X)$, the following are equivalent:
\begin{enumerate}
\item $\operatorname{EP}_X(\theta)$;
\item the normalized space
\[
 \PSH_X^{\sup}(\theta)
   =\{\phi\in\PSH_X(\theta)\mid\sup\phi=0\}
\]
is compact for the topology of pointwise convergence on
$X^{\divis}$;
\item every bounded-above increasing net in $\PSH_X(\theta)$ converges
      pointwise on $X^{\divis}$ to the divisorial trace of a unique
      $\theta$-psh function.
\end{enumerate}
In case (3), the resulting psh function is the usc regularization of the raw
pointwise supremum.
\end{theorem}

\begin{proof}
Assume (1) and take a net in the normalized space.  The BJ pointwise bounds
put every divisorial coordinate in a compact interval.  Tychonoff gives a
pointwise convergent subnet.  For each index $i$, take the regularized
supremum of the tail.  By (1) this is psh; the tail envelopes decrease, and
their limit is psh because its value at $v_{X,\triv}$ is zero.  The
nonnegligibility of divisorial points identifies its trace with the
coordinatewise cluster point.  Thus (2) holds.

Assume (2).  Normalize a bounded-above increasing net by subtracting its
supremum.  A subnet converges in the normalized compact space, while the
subtracted real constants converge.  Monotonicity forces the entire net to
have the same divisorial limit, proving (3).

Finally assume (3), and let $(\phi_i)$ be a nonempty bounded-above family.
For each finite nonempty $J\subset I$, put
$\psi_J=\max_{i\in J}\phi_i$.  The $\psi_J$ form an increasing
bounded-above net, so (3) gives $\psi\in\PSH_X(\theta)$.  On divisorial
points,
\[
             \psi=\sup_J\psi_J=\sup_i\phi_i
                   =\left(\sup_i\phi_i\right)^*,
\]
where the last equality is divisorial nonnegligibility.  Since $\psi$ is
psh and the regularized supremum is usc, divisorial comparison gives
$\psi\leq(\sup_i\phi_i)^*$.  Conversely, comparison between the psh
functions $\phi_i$ and $\psi$ gives $\phi_i\leq\psi$ on all of
$X^{\An}$.  Taking the raw supremum and then its usc regularization gives
$(\sup_i\phi_i)^*\leq\psi$.  This proves (1).
This is \cite[Theorem~5.11]{BJ22}.
\end{proof}

\begin{warning}{Why a sequence is insufficient}
\label{warn:sequence-insufficient-envelope}
The directed set of finite subsets of an arbitrary index set need not have a
countable cofinal subset.  Replacing the canonical net $(\psi_J)$ by a
sequence therefore changes the proof unless a countable subfamily with the
same supremum is separately produced.  No such reduction is part of the
general definition.
\end{warning}

\section{Obstacle formulations}

\begin{theorem}\label{thm:general-obstacle-criterion}
Assume $\PSH_X(\theta)\ne\varnothing$.  Then
$\operatorname{EP}_X(\theta)$ is equivalent to each of the following:
\begin{enumerate}
\item for every extended-real obstacle $h$, one has exactly one of
\[
 P_\theta(h)\equiv-\infty,\qquad
 P_\theta(h)^*\equiv+\infty,\qquad
 P_\theta(h)^*\in\PSH_X(\theta);
\tag{36.4}\label{eq:envelope-trichotomy}
\]
\item $P_\theta(h)\in\PSH_X(\theta)$ for every
      $h\in\PL(X)$.
\end{enumerate}
\par
\end{theorem}

\begin{proof}
Under EP, the competitors in \eqref{eq:bj-obstacle-envelope} either are
uniformly bounded above, in which case their regularized supremum is psh, or
their suprema tend to $+\infty$.  The Alexander--Taylor bound then forces
the regularization to be $+\infty$ on all divisorial points and hence
identically $+\infty$.  This proves \eqref{eq:envelope-trichotomy}.

For a PL obstacle, boundedness and the existence of one psh function rule out
both infinite alternatives.  Since the obstacle is continuous,
$P_\theta(h)^*\leq h$, so the regularization is itself a competitor and
equals the raw envelope.

Conversely, let $(\phi_i)$ be bounded above and put
$h=(\sup_i\phi_i)^*$.  Since $h$ is upper semicontinuous on the compact
Hausdorff space $X^{\An}$, it is the decreasing limit of a net of
continuous majorants.  Uniform density \eqref{eq:pl-stone-weierstrass},
together with closure of $\PL(X)$ under finite minima, refines this to PL
majorants $h_j\searrow h$.  Then $P_\theta(h_j)$ is psh and
\[
          \sup_i\phi_i\leq P_\theta(h_j)\leq h_j.
\]
Monotonicity of the obstacle envelope makes $P_\theta(h_j)$ decreasing.
Its limit is psh or identically $-\infty$ by
\Cref{thm:bj-basic-closure}\,(4), and the lower bound supplied by one fixed
$\phi_i$ excludes the latter.  Passing to the limit in the displayed
squeeze gives $P_\theta(h_j)\searrow h$, so $h$ is psh.
This is the criterion of \cite[Lemma~1.3]{BJ22b}.
\end{proof}

\begin{proposition}\label{prop:ch36-continuity-of-envelopes}
For an ample class $\omega$, the envelope property is equivalent to
continuity of envelopes:
\[
 h\in C^0(X^{\An})\quad\Longrightarrow\quad
                     P_\omega(h)\in C^0(X^{\An})\cap\PSH_X(\omega).
\tag{36.5}\label{eq:continuity-of-envelopes}
\]
\end{proposition}

\begin{proof}
Assume EP and let $h$ be continuous.  The general obstacle criterion makes
$P_\omega(h)$ psh and hence upper semicontinuous.  If
$\psi\in\PSH_X(\omega)$ satisfies $\psi\leq h$, choose a decreasing net
$\psi_i\in\cH_{\mathrm{dom}}(\omega)$ converging to $\psi$.  Compactness,
upper semicontinuity of $\psi_i$, and lower semicontinuity of $h$ give,
for each $\varepsilon>0$, an index $i$ with
$\psi\leq\psi_i\leq h+\varepsilon$.  Hence
$\psi-\varepsilon\leq\psi_i-\varepsilon\leq h$, with
$\psi_i-\varepsilon$ continuous and psh.  Taking the supremum over
$\psi$ and letting $\varepsilon\downarrow0$ shows that the envelope is
lower semicontinuous; it is therefore continuous.

Conversely, let $(\phi_i)$ be bounded above and set
$h=(\sup_i\phi_i)^*$.  Choose continuous majorants $h_j\searrow h$.
By assumption, each $P_\omega(h_j)$ is continuous and psh, while
$\sup_i\phi_i\leq P_\omega(h_j)\leq h_j$.  The same decreasing-limit
argument as above gives $P_\omega(h_j)\searrow h$, so $h$ is psh.
This is \cite[Lemma~5.17]{BJ22}.
\end{proof}


\section{Nets, sequences, and increasing families}
\label{sec:ch36-nets}

This section makes precise the third and fourth items of
\Cref{sec:ch36-orientation}.  The envelope theorem concerns arbitrary
families, and their canonical directed replacement is indexed by finite
subsets.  The indexing obstruction below is all that is needed here.

\subsection{Why a sequence will not do}

\begin{proposition}
\label{prop:ch36-no-sequences}
For a nonempty set $A$, let $\mathcal F(A)$ be the directed set of nonempty finite
subsets of $A$, ordered by inclusion.  Then $\mathcal F(A)$ has a
countable cofinal subset if and only if $A$ is countable.  Consequently,
when a family $(\phi_\alpha)_{\alpha\in A}$ is replaced by its finite
maxima, that formal reduction cannot in general be replaced by a sequence if
$A$ is uncountable.
\end{proposition}

\begin{proof}
If $A$ is finite, its largest element in $\mathcal F(A)$ is $A$ itself.  If
$A=\{a_1,a_2,\dots\}$ is countably infinite, the sets
$\{a_1,\dots,a_n\}$ form a cofinal sequence.  Conversely, if
$(S_n)_{n\geq1}\subset\mathcal F(A)$ is cofinal, then every singleton
$\{a\}$ is contained in some $S_n$, so
$A=\bigcup_{n\geq1}S_n$ is countable.
\end{proof}

This proposition concerns the canonical finite-maximum construction, not the
existence of a specially chosen countable subfamily with the same supremum.
Accordingly, later arguments retain nets unless the data are already
countably indexed or a separate sequential reduction is proved.

\subsection{Directedness is free; upper semicontinuity is not}

\begin{lemma}
\label{lem:ch36-directed-reduction}
Let $(\phi_\alpha)_{\alpha\in A}\subset\PSH_X(\theta)$ be nonempty and
bounded above.  Let $F$ be the set of nonempty finite subsets of $A$,
directed by inclusion, and for $S\in F$ put
$\phi_S:=\max_{\alpha\in S}\phi_\alpha$.
Then $(\phi_S)_{S\in F}$ is an increasing net in $\PSH_X(\theta)$ with
$\sup_S\phi_S=\sup_\alpha\phi_\alpha$, and hence
\[
 \Bigl(\sup_{\alpha}\phi_\alpha\Bigr)^*=\Bigl(\sup_{S}\phi_S\Bigr)^* .
\]
\end{lemma}

\begin{proof}
Each $\phi_S$ is $\theta$-psh by \Cref{thm:bj-basic-closure}\,(2) and
induction on $\#S$; the net is increasing because $S\subset S'$ implies
$\phi_S\le\phi_{S'}$; and the two suprema agree pointwise since every
$\phi_\alpha$ equals $\phi_{\{\alpha\}}$.  Equality of the
regularizations follows because they are computed from equal functions.
\end{proof}

So the passage from an arbitrary family to a directed one costs nothing.
What is not free is the passage from the raw supremum to its regularization,
and the following example shows the two genuinely differ.

\begin{example}[Why a regularization step is unavoidable]
\label{ex:ch36-sup-not-usc}
The passage from $\sup_i\phi_i$ to $(\sup_i\phi_i)^*$ in
\eqref{eq:bj-envelope-property} is not cosmetic: a supremum of upper
semicontinuous functions need not be upper semicontinuous, already on a
segment.  On the skeleton $[0,1]$ put $f_j(t):=\min\{1,jt\}$ for
$j\in\mathbf Z_{>0}$.  Each $f_j$ is continuous, the family is
increasing, and
\[
 \sup_j f_j(t)=\begin{cases}0,&t=0,\\ 1,&t>0,\end{cases}
\]
which fails upper semicontinuity at $t=0$, where the value is $0$ while
$\limsup_{t\to0^+}=1$.  Its regularization is the constant $1$.

The example is elementary and is not itself about psh functions; it isolates
the mechanism.  An increasing family can leave a deficit on a closed set,
and the regularization fills it in.  The envelope property asserts that the
regularized function is still in $\PSH_X(\theta)$; it does not say that one
member of the original family attains the new value.  This is why
\Cref{def:bj-envelope-property} regularizes first and asks for membership
second, and why no formal argument from
\Cref{thm:bj-basic-closure} can prove it: that theorem's only limit
operation is decreasing, and decreasing limits of usc functions are usc for
free.
\end{example}

\begin{warning}{The obstacle envelope is not automatically an operator}
Before the envelope theorem, $P_\theta$ of
\Cref{def:bj-obstacle-envelope} is a map from obstacles to
\emph{functions}, not to $\PSH_X(\theta)$.  The formal identity
$P_\theta(P_\theta(h))=P_\theta(h)$ in
\Cref{lem:obstacle-envelope-formal}\,(3) compares two competitor sets and
does not make $P_\theta(h)$ psh.  Writing $P_\theta$ as a self-map of the
psh class presupposes exactly what is at issue.  Every occurrence before
\Cref{chap:bj-envelope-statement} in this book is therefore accompanied by
an explicit statement of which of the three alternatives after
\eqref{eq:bj-obstacle-envelope} is being assumed or proved.
\end{warning}

\section{Necessary geometry: unibranchness}

Recall that an integral variety is unibranch if its normalization
$\nu:X^\nu\to X$ is a homeomorphism on underlying Zariski spaces.
Normal varieties are unibranch.  Over $\mathbf C$, this is equivalent to
local analytic irreducibility.

\begin{theorem}\label{thm:ep-implies-unibranch}
If $\operatorname{EP}_X(\omega)$ holds for one ample class
$\omega$, then $X$ is unibranch.
\end{theorem}

\begin{proof}
Let $\nu:X^\nu\to X$ be the normalization.  For
$\omega_m=(1+m^{-1})\omega$, scaling shows that $\omega_m$ has the
envelope property and $\omega_m>\omega$.  The birational pullback theorem
\cite[Lemma~5.13]{BJ22}, whose proof uses
\Cref{thm:bj-envelope-equivalences}, therefore makes
$\nu^*:\PSH_X(\omega)\to\PSH_{X^\nu}(\nu^*\omega)$ surjective.  In
particular, every function in
$\cH_{\mathrm{dom}}(\nu^*\omega)$ descends to $X^{\An}$.  These
functions span $\PL(X^\nu)$, which separates points of
$(X^\nu)^{\An}$; hence $\nu^{\An}$ is injective.  Non-Archimedean GAGA
then makes the finite normalization injective on scheme points.  Since it is
surjective, it is a homeomorphism, and $X$ is unibranch.  This is
\cite[Theorem~5.12]{BJ22}.
\end{proof}

\begin{example}[The node and the cusp]
\label{ex:node-cusp-envelope}
For a nodal curve, normalization has two points above the node.  Their
proper-support semivaluations are identified in $X^{\An}$, while PL
functions upstairs can separate them.  If EP held in an ample class, the
descent argument above would force every such separator to descend, a
contradiction.  A cusp is unibranch: normalization is bijective on points,
though not an isomorphism.  Thus the node is excluded by necessity, while
the cusp belongs to the class predicted by the historical conjecture.
\end{example}

\section{Published smooth cases}

The first general positive result was:

\begin{theorem}\label{thm:bj22-nef-envelope}
Assume that $X$ is smooth and connected, and that either
$\operatorname{char}k=0$ or $\dim X\leq2$.  Then every nef numerical
class $\theta$ has the envelope property.
\end{theorem}

\begin{proof}
This is \cite[Theorem~5.20]{BJ22}; we reproduce its reduction and the
uniform estimate.  By \cite[Lemma~5.9]{BJ22}, it is enough to treat
$\theta=c_1(L)$ for an ample $\mathbf Q$-line bundle $L$: indeed, a
nef real class is the limit of ample rational classes $\theta_j\geq\theta$,
and the decreasing-limit axiom passes the envelope property to the limit.
By \Cref{prop:ch36-continuity-of-envelopes}, it remains
to prove that $P_L(h)$ is continuous for $h\in C^0(X^{\An})$.  Since
\[
 \lVert P_L(h)-P_L(h')\rVert_\infty
       \leq \lVert h-h'\rVert_\infty,
\]
uniform density of $\PL(X)$ reduces this to $h\in\PL(X)$.

The test-configuration realization
\cite[Theorem~2.31]{BJ22} writes $h=\phi_{\cL}$ for an integrally closed
test configuration $(\cX,\cL)$ of $(X,L)$.  After clearing
denominators, let $\mathfrak a_m$ be the base ideal of $m\cL$, and put
\[
 u_m:=\phi_{\cL}+m^{-1}\phi_{\mathfrak a_m}.
\]
The inclusions
$\mathfrak a_m\mathfrak a_{m'}\subset\mathfrak a_{m+m'}$ and the
valuative description of sections give
\[
                 u_m\longrightarrow P_L(\phi_{\cL})
                 \quad\text{pointwise}.
\tag{36.5a}\label{eq:bj22-nef-base-ideal-limit}
\]
This is \cite[Lemma~5.22]{BJ22}.  Resolving the total space, we may assume
that $\cX$ is smooth with SNC central fibre; in the positive-characteristic
case this is where $\dim\cX\leq3$ is used.

Suppose first that $\operatorname{char}k=0$, and let
$\mathfrak b_m$ be the asymptotic multiplier ideal of the graded sequence
$\mathfrak a_\bullet$.  The two foundational ideal inputs are
\[
 \mathfrak a_m\subset\mathfrak b_m,
 \qquad
 \mathfrak b_{m\ell}\subset\mathfrak b_m^{\,\ell}.
\tag{36.5b}\label{eq:bj22-nef-multiplier-inclusions}
\]
They imply, after letting $\ell\to\infty$ in
\eqref{eq:bj22-nef-base-ideal-limit},
\[
 P_L(\phi_{\cL})
 \leq \phi_{\cL}+m^{-1}\phi_{\mathfrak b_m}.
\tag{36.5c}\label{eq:bj22-nef-multiplier-upper}
\]
Uniform global generation supplies an effective vertical
$\mathbf Q$-divisor $E$, an ample
$\mathcal A:=\cL-E$, an integer $m_0$, and a constant $C$, independent of
$m$, such that
\[
 \phi_{\mathfrak b_m}\leq
 \phi_{\mathfrak a_{m+m_0}}+C.
\]
Combining this with \eqref{eq:bj22-nef-multiplier-upper} gives the uniform
squeeze
\[
 u_m\leq P_L(\phi_{\cL})
 \leq \left(1+\frac{m_0}{m}\right)u_{m+m_0}
      -\frac{m_0}{m}\phi_{\cL}+\frac Cm.
\tag{36.5d}\label{eq:bj22-nef-uniform-squeeze}
\]
Together with the superadditivity used in
\eqref{eq:bj22-nef-base-ideal-limit},
\eqref{eq:bj22-nef-uniform-squeeze} makes that convergence uniform.  Hence
the envelope is continuous.  When $\operatorname{char}k>0$, the same inclusions and
uniform-generation estimate hold for the test ideals used in the proof of
\cite[Theorem~5.20]{BJ22}; the precise imported test-ideal inputs are
subadditivity and uniform generation in
\cite[\S4 and Theorem~4.6]{GJKM19}.  The identical squeeze completes the
proof.
\end{proof}

The proof adapts the model-envelope regularity argument from the discretely
valued setting.  In characteristic zero it uses multiplier ideals; in the
low-dimensional positive-characteristic range it uses the available test
ideal technology.  Smoothness and the stated characteristic/dimension
alternative are part of the theorem, not expository conveniences.

The 2022 preprint addendum, in its revised 2023 form and published in 2024,
removed the nefness restriction:

For a big $\mathbf Q$-line bundle $L$ and a bounded function $h$, set
\[
 Q_L(h):=\sup\{\psi\in\cH^{\mathrm{gf}}_{\mathbf R}(L)\mid\psi\leq h\}.
\tag{36.6}\label{eq:fs-envelope}
\]
This is the Fubini--Study envelope used in the proof below; it is distinct
from the all-psh obstacle envelope $P_L$.

\begin{theorem}\label{thm:bj-addendum-all-classes}
Assume that $X$ is smooth and that either
$\operatorname{char}k=0$ or $\dim X\leq2$.  Then every
$\theta\in N^1(X)$ has the envelope property.
\end{theorem}

\begin{proof}
This is Theorem~A of \cite{BJ22b}.  If $\theta$ is not pseudoeffective,
then $\PSH_X(\theta)=\varnothing$, so the assertion is vacuous.  Otherwise
choose big $\mathbf Q$-line bundles $L_i$ with
$c_1(L_i)\geq\theta$ and $c_1(L_i)\to\theta$.  The same limit argument
as in \cite[Lemma~5.9]{BJ22} reduces the proof to
$\theta=c_1(L)$ with $L$ big.  By
\Cref{thm:general-obstacle-criterion}, it is enough to show that
$P_L(h)$ is $L$-psh for every $h\in\PL(X)$.

Realize $h=\phi_{\cL}$ on an integrally closed test configuration
$(\cX,\cL)$, and clear denominators.  If $\mathfrak a_m$ is the base
ideal of $m\cL$, then \cite[Lemmas~1.6--1.7]{BJ22b} gives
\[
 q_m:=\phi_{\cL}+m^{-1}\phi_{\mathfrak a_m}
       \nearrow Q_L(\phi_{\cL})
\tag{36.7}\label{eq:addendum-q-limit}
\]
along the divisibility order.  The membership test behind this equality is
exact: a Laurent section
$s=\sum_\lambda s_\lambda\varpi^{-\lambda}$ belongs to
$H^0(\cX,m\cL)$ if and only if
\[
 \max_\lambda\{\log|s_\lambda|+\lambda\}
       \leq\phi_{m\cL}
\]
on $X^{\An}$.  Thus every GF competitor below $\phi_{\cL}$ occurs in
one of the base-ideal stages, proving \eqref{eq:addendum-q-limit} rather
than only one inequality.  The index set consists of sufficiently divisible
positive integers ordered by divisibility, not by their usual order.

After resolving $\cX$, let $\mathfrak b_m$ denote the asymptotic
multiplier ideals of $\mathfrak a_\bullet$ in characteristic zero, and
the corresponding asymptotic test ideals in positive characteristic.  The
inclusions
\[
 \mathfrak a_m\subset\mathfrak b_m,
 \qquad
 \mathfrak b_{m\ell}\subset\mathfrak b_m^{\,\ell}
\tag{36.8}\label{eq:addendum-ideal-inclusions}
\]
give
\[
 Q_L(\phi_{\cL})\leq
 \psi_m:=\phi_{\cL}+m^{-1}\phi_{\mathfrak b_m}.
\tag{36.9}\label{eq:addendum-Q-upper}
\]
Uniform global generation provides a $\mathbf G_m$-equivariant ample line
bundle $\mathcal A$ on $\cX$ such that
$\cO_{\cX}(m\cL+\mathcal A)\otimes\mathfrak b_m$ is globally generated for all
admissible $m$.  Adding a multiple of $\cX_0$, arrange
$\phi_{\mathcal A}\geq0$.  Then
\[
 \psi'_m:=\psi_m+m^{-1}\phi_{\mathcal A}
     \in\cH^{\mathrm{gf}}_{\mathbf Q}(L+m^{-1}A)
\]
decreases in the divisibility order.  Hence
$\psi:=\inf_m\psi'_m$ is $L$-psh or identically $-\infty$.  But
\eqref{eq:addendum-Q-upper} and global generation give
\[
 Q_L(\phi_{\cL})\leq\psi'_m
       \leq\phi_{\cL}+m^{-1}\phi_{\mathcal A},
\]
so
\[
                 Q_L(\phi_{\cL})\leq\psi\leq\phi_{\cL};
\tag{36.10}\label{eq:addendum-psi-squeeze}
\]
in particular, $\psi\not\equiv-\infty$.

It remains to identify the limit.  By \cite[Lemma~1.5]{BJ22b},
$P_L(\phi_{\cL})=Q_L(\phi_{\cL})$ over the complement of the augmented
base locus.  Every divisorial valuation has generic support, hence lies in
that analytic open set.  Thus any $\tau\in\PSH_X(L)$ with
$\tau\leq\phi_{\cL}$ satisfies $\tau\leq\psi$ on $X^{\divis}$, and divisorial comparison
\cite[Theorem~4.22]{BJ22} gives the inequality on all of $X^{\An}$.
Taking the supremum over $\tau$ yields
$P_L(\phi_{\cL})\leq\psi$.  Conversely,
\eqref{eq:addendum-psi-squeeze} makes $\psi$ a competitor, so
$\psi\leq P_L(\phi_{\cL})$.  Therefore the obstacle envelope is psh,
and \Cref{thm:general-obstacle-criterion} proves the theorem.  The only
imported characteristic-dependent core is precisely the multiplier-ideal or
test-ideal package \eqref{eq:addendum-ideal-inclusions} with uniform global
generation, anchored in \cite[\S2.1]{BJ22b}, with the
positive-characteristic subadditivity and uniform-generation inputs given by
\cite[\S4 and Theorem~4.6]{GJKM19}.
\end{proof}

The same addendum proved a separate energy statement on singular varieties.
If $\omega$ is ample and $(X,\omega)$ has the weak envelope property
(there is a birational model with an ample dominating class having EP), then
for every bounded $h$,
\[
E_\omega(h)=E_\omega(P_\omega(h))
             =E_\omega(P_\omega(h)^*).
\tag{36.11}\label{eq:addendum-energy-envelope}
\]
Here $E_\omega$ is the trivially valued Boucksom--Jonsson energy of
\cite{BJ22,BJ22b}; no discretely valued BFJ model-form pairing is being
used.
This does not assert that $P_\omega(h)^*$ is psh when EP fails.  Equality
of energy is therefore strictly weaker than envelope regularity.
This is \cite[Theorem~B]{BJ22b}.

\section{The historical gap and the proof map ahead}

At the end of the 2022 theory and its addendum, the exact unresolved
statement was:

\begin{quote}
\emph{If $X$ is an integral projective unibranch variety over an
algebraically closed trivially valued field, then
$\operatorname{EP}_X(\theta)$ holds for every
$\theta\in N^1(X)_{\mathbf R}$.}
\end{quote}

This was Conjecture~5.14 of \cite{BJ22}.  It asked simultaneously for
singular unibranch varieties, arbitrary dimension, arbitrary characteristic,
and arbitrary real numerical classes.  The smooth theorem above did not
settle descent through normalization or alterations.  Divisorial equality
alone was also insufficient: a GF function may be $-\infty$ on
proper-support points.

The next chapters prove the statement, without using it as an input.  The
dependency map is:
\[
\begin{array}{c|l}
\text{Chapter}&\text{new mechanism}\\ \hline
37&\text{final statement, sharpness, and proof ledger}\\
38&\text{Cartier b-divisor realization independent of EP}\\
39&\text{birational and homogeneous-carrier reflection}\\
40&\text{smooth case in arbitrary characteristic}\\
41&\text{finite normal descent by determinant norms}\\
42&\text{alterations and the normal case}\\
43&\text{universal-homeomorphic descent in characteristic $p$}\\
44&\text{subintegral descent in characteristic zero}\\
45&\text{assembly and consequences}.
\end{array}
\]

The logical firewall is important: none of
\Crefrange{thm:bj-envelope-equivalences}{thm:bj-addendum-all-classes}
contains the later general theorem.  This chapter records the historical
inputs; \Crefrange{chap:bj-envelope-statement}{chap:bj-envelope-consequences}
supplies the missing descent and realization arguments.

\section{Source note}

\Cref{def:bj-envelope-property,def:bj-obstacle-envelope},
\Cref{thm:bj-envelope-equivalences,prop:ch36-continuity-of-envelopes,%
thm:ep-implies-unibranch,thm:bj22-nef-envelope}, and the historical
conjecture are respectively Definition~5.8, Theorem~5.11, Lemma~5.17,
Theorems~5.12 and~5.20, and Conjecture~5.14 of \cite{BJ22}.  The general
obstacle criterion is \cite[Lemma~1.3]{BJ22b}; the all-class theorem and the
energy identity are Theorems~A and~B there, while
\eqref{eq:addendum-q-limit} is its Lemma~1.6.  The positive-characteristic
ideal inputs are the subadditivity statement in
\cite[\S4, Remark~4.4(3)]{GJKM19} and uniform generation
\cite[Theorem~4.6]{GJKM19}.  The conjecture is stated here only at its
historical location; the complete solution is the subject of the following
nine chapters.
